\subsection{Nefakultní budovy}
Kromě budov, ve kterých probíhá výuka těch podřadnějších předmětů, jako je
matematická analýza, navštívíte i budovy jiné. Nejdůležitější z nich je SCUK
Hostivař, kde probíhá většina tělocviků. Dalšími budovami jsou například
Karolinum – historické sídlo univerzity, které navštívíte při imatrikulaci – a
nebo loděnice v Malé Chuchli, kde probíhá výuka kanoistiky.


\subsubsection{SCUK (Sportovní centrum UK) Hostivař}
\budovaKratce{Bruslařská 10, 102 00 Praha 10 - Hostivař}
{Tělocvikáři}
V prvním patře (to je to, do kterého vejdete hlavním vchodem) této budovy sídlí
katedry tělesné výchovy (KTV) několika fakult, naše KTV je samozřejmě ta úplně vzadu.

Pokud hledáte šatny, u většiny sportů půjdete správně, když za vchodem sejdete
po schodech a dáte se do chodby vlevo (dámské jsou pak vpravo, pánské vlevo). Na
bazén si oproti ISICu vyzvedněte visací zámeček u okýnka pod schody a jděte do
šaten za okýnkem. Na ostatní sporty byste měli mít zámek vlastní.

Tato budova, nebo spíše její umístění, se stane důvodem, proč budou někteří
z~vás nenávidět tělesnou výchovu. Budova se totiž nachází na opačném konci Prahy
než zbytek matfyzáckého světa.


\subsubsection{Karolinum (rektorát UK)}
\budovaKratce{Ovocný trh 5, 116 36 Praha 1}
{Byrokraté}
Rektorát fakultě pochopitelně nepatří. Je to něco jako náš děkanát, ale „o patro
výš“. Ve staroslavném \textit{Karolinu} se konají dva obřady ohraničující
studium: \textit{imatrikulace} (tu může zažít každý) a \textit{promoce} (tu
jenom ti vytrvalí).


\subsubsection{Loděnice Malá Chuchle}
\budovaKratce{Strakonická Praha 5 - Malá Chuchle}
{Tělocvikáři}
Na ty, kteří si zapíší kanoistiku, čeká první měsíc zimního semestru a poslední
měsíc(e) semestru letního dojíždění do Malé Chuchle (v zimě se s kajaky nacpete
mezi plavce do bazénu v Hostivaři).